% Options for packages loaded elsewhere
% Options for packages loaded elsewhere
\PassOptionsToPackage{unicode}{hyperref}
\PassOptionsToPackage{hyphens}{url}
\PassOptionsToPackage{dvipsnames,svgnames,x11names}{xcolor}
%
\documentclass[
  english,
  10pt,
]{article}
\usepackage{xcolor}
\usepackage[margin=1in]{geometry}
\usepackage{amsmath,amssymb}
\setcounter{secnumdepth}{-\maxdimen} % remove section numbering
\usepackage{iftex}
\ifPDFTeX
  \usepackage[T1]{fontenc}
  \usepackage[utf8]{inputenc}
  \usepackage{textcomp} % provide euro and other symbols
\else % if luatex or xetex
  \usepackage{unicode-math} % this also loads fontspec
  \defaultfontfeatures{Scale=MatchLowercase}
  \defaultfontfeatures[\rmfamily]{Ligatures=TeX,Scale=1}
\fi
\usepackage{lmodern}
\ifPDFTeX\else
  % xetex/luatex font selection
  \setmainfont[]{Latin Modern Roman}
  \setsansfont[]{Latin Modern Sans}
  \setmonofont[]{Latin Modern Mono}
\fi
% Use upquote if available, for straight quotes in verbatim environments
\IfFileExists{upquote.sty}{\usepackage{upquote}}{}
\IfFileExists{microtype.sty}{% use microtype if available
  \usepackage[]{microtype}
  \UseMicrotypeSet[protrusion]{basicmath} % disable protrusion for tt fonts
}{}
\makeatletter
\@ifundefined{KOMAClassName}{% if non-KOMA class
  \IfFileExists{parskip.sty}{%
    \usepackage{parskip}
  }{% else
    \setlength{\parindent}{0pt}
    \setlength{\parskip}{6pt plus 2pt minus 1pt}}
}{% if KOMA class
  \KOMAoptions{parskip=half}}
\makeatother
% Make \paragraph and \subparagraph free-standing
\makeatletter
\ifx\paragraph\undefined\else
  \let\oldparagraph\paragraph
  \renewcommand{\paragraph}{
    \@ifstar
      \xxxParagraphStar
      \xxxParagraphNoStar
  }
  \newcommand{\xxxParagraphStar}[1]{\oldparagraph*{#1}\mbox{}}
  \newcommand{\xxxParagraphNoStar}[1]{\oldparagraph{#1}\mbox{}}
\fi
\ifx\subparagraph\undefined\else
  \let\oldsubparagraph\subparagraph
  \renewcommand{\subparagraph}{
    \@ifstar
      \xxxSubParagraphStar
      \xxxSubParagraphNoStar
  }
  \newcommand{\xxxSubParagraphStar}[1]{\oldsubparagraph*{#1}\mbox{}}
  \newcommand{\xxxSubParagraphNoStar}[1]{\oldsubparagraph{#1}\mbox{}}
\fi
\makeatother

\usepackage{color}
\usepackage{fancyvrb}
\newcommand{\VerbBar}{|}
\newcommand{\VERB}{\Verb[commandchars=\\\{\}]}
\DefineVerbatimEnvironment{Highlighting}{Verbatim}{commandchars=\\\{\}}
% Add ',fontsize=\small' for more characters per line
\usepackage{framed}
\definecolor{shadecolor}{RGB}{241,243,245}
\newenvironment{Shaded}{\begin{snugshade}}{\end{snugshade}}
\newcommand{\AlertTok}[1]{\textcolor[rgb]{0.68,0.00,0.00}{#1}}
\newcommand{\AnnotationTok}[1]{\textcolor[rgb]{0.37,0.37,0.37}{#1}}
\newcommand{\AttributeTok}[1]{\textcolor[rgb]{0.40,0.45,0.13}{#1}}
\newcommand{\BaseNTok}[1]{\textcolor[rgb]{0.68,0.00,0.00}{#1}}
\newcommand{\BuiltInTok}[1]{\textcolor[rgb]{0.00,0.23,0.31}{#1}}
\newcommand{\CharTok}[1]{\textcolor[rgb]{0.13,0.47,0.30}{#1}}
\newcommand{\CommentTok}[1]{\textcolor[rgb]{0.37,0.37,0.37}{#1}}
\newcommand{\CommentVarTok}[1]{\textcolor[rgb]{0.37,0.37,0.37}{\textit{#1}}}
\newcommand{\ConstantTok}[1]{\textcolor[rgb]{0.56,0.35,0.01}{#1}}
\newcommand{\ControlFlowTok}[1]{\textcolor[rgb]{0.00,0.23,0.31}{\textbf{#1}}}
\newcommand{\DataTypeTok}[1]{\textcolor[rgb]{0.68,0.00,0.00}{#1}}
\newcommand{\DecValTok}[1]{\textcolor[rgb]{0.68,0.00,0.00}{#1}}
\newcommand{\DocumentationTok}[1]{\textcolor[rgb]{0.37,0.37,0.37}{\textit{#1}}}
\newcommand{\ErrorTok}[1]{\textcolor[rgb]{0.68,0.00,0.00}{#1}}
\newcommand{\ExtensionTok}[1]{\textcolor[rgb]{0.00,0.23,0.31}{#1}}
\newcommand{\FloatTok}[1]{\textcolor[rgb]{0.68,0.00,0.00}{#1}}
\newcommand{\FunctionTok}[1]{\textcolor[rgb]{0.28,0.35,0.67}{#1}}
\newcommand{\ImportTok}[1]{\textcolor[rgb]{0.00,0.46,0.62}{#1}}
\newcommand{\InformationTok}[1]{\textcolor[rgb]{0.37,0.37,0.37}{#1}}
\newcommand{\KeywordTok}[1]{\textcolor[rgb]{0.00,0.23,0.31}{\textbf{#1}}}
\newcommand{\NormalTok}[1]{\textcolor[rgb]{0.00,0.23,0.31}{#1}}
\newcommand{\OperatorTok}[1]{\textcolor[rgb]{0.37,0.37,0.37}{#1}}
\newcommand{\OtherTok}[1]{\textcolor[rgb]{0.00,0.23,0.31}{#1}}
\newcommand{\PreprocessorTok}[1]{\textcolor[rgb]{0.68,0.00,0.00}{#1}}
\newcommand{\RegionMarkerTok}[1]{\textcolor[rgb]{0.00,0.23,0.31}{#1}}
\newcommand{\SpecialCharTok}[1]{\textcolor[rgb]{0.37,0.37,0.37}{#1}}
\newcommand{\SpecialStringTok}[1]{\textcolor[rgb]{0.13,0.47,0.30}{#1}}
\newcommand{\StringTok}[1]{\textcolor[rgb]{0.13,0.47,0.30}{#1}}
\newcommand{\VariableTok}[1]{\textcolor[rgb]{0.07,0.07,0.07}{#1}}
\newcommand{\VerbatimStringTok}[1]{\textcolor[rgb]{0.13,0.47,0.30}{#1}}
\newcommand{\WarningTok}[1]{\textcolor[rgb]{0.37,0.37,0.37}{\textit{#1}}}

\usepackage{longtable,booktabs,array}
\newcounter{none} % for unnumbered tables
\usepackage{calc} % for calculating minipage widths
% Correct order of tables after \paragraph or \subparagraph
\usepackage{etoolbox}
\makeatletter
\patchcmd\longtable{\par}{\if@noskipsec\mbox{}\fi\par}{}{}
\makeatother
% Allow footnotes in longtable head/foot
\IfFileExists{footnotehyper.sty}{\usepackage{footnotehyper}}{\usepackage{footnote}}
\makesavenoteenv{longtable}
\usepackage{graphicx}
\makeatletter
\newsavebox\pandoc@box
\newcommand*\pandocbounded[1]{% scales image to fit in text height/width
  \sbox\pandoc@box{#1}%
  \Gscale@div\@tempa{\textheight}{\dimexpr\ht\pandoc@box+\dp\pandoc@box\relax}%
  \Gscale@div\@tempb{\linewidth}{\wd\pandoc@box}%
  \ifdim\@tempb\p@<\@tempa\p@\let\@tempa\@tempb\fi% select the smaller of both
  \ifdim\@tempa\p@<\p@\scalebox{\@tempa}{\usebox\pandoc@box}%
  \else\usebox{\pandoc@box}%
  \fi%
}
% Set default figure placement to htbp
\def\fps@figure{htbp}
\makeatother



\ifLuaTeX
\usepackage[bidi=basic,shorthands=off]{babel}
\else
\usepackage[bidi=default,shorthands=off]{babel}
\fi
\ifPDFTeX
\else
\babelfont{rm}[]{Latin Modern Roman}
\fi
\ifLuaTeX
  \usepackage{selnolig} % disable illegal ligatures
\fi


\setlength{\emergencystretch}{3em} % prevent overfull lines

\providecommand{\tightlist}{%
  \setlength{\itemsep}{0pt}\setlength{\parskip}{0pt}}






\usepackage{setspace}
\onehalfspacing
\usepackage{booktabs}
\usepackage{longtable}
\usepackage{array}
\usepackage{xcolor}
\makeatletter
\@ifpackageloaded{tcolorbox}{}{\usepackage[skins,breakable]{tcolorbox}}
\@ifpackageloaded{fontawesome5}{}{\usepackage{fontawesome5}}
\definecolor{quarto-callout-color}{HTML}{909090}
\definecolor{quarto-callout-note-color}{HTML}{0758E5}
\definecolor{quarto-callout-important-color}{HTML}{CC1914}
\definecolor{quarto-callout-warning-color}{HTML}{EB9113}
\definecolor{quarto-callout-tip-color}{HTML}{00A047}
\definecolor{quarto-callout-caution-color}{HTML}{FC5300}
\definecolor{quarto-callout-color-frame}{HTML}{acacac}
\definecolor{quarto-callout-note-color-frame}{HTML}{4582ec}
\definecolor{quarto-callout-important-color-frame}{HTML}{d9534f}
\definecolor{quarto-callout-warning-color-frame}{HTML}{f0ad4e}
\definecolor{quarto-callout-tip-color-frame}{HTML}{02b875}
\definecolor{quarto-callout-caution-color-frame}{HTML}{fd7e14}
\makeatother
\makeatletter
\@ifpackageloaded{caption}{}{\usepackage{caption}}
\AtBeginDocument{%
\ifdefined\contentsname
  \renewcommand*\contentsname{Table of contents}
\else
  \newcommand\contentsname{Table of contents}
\fi
\ifdefined\listfigurename
  \renewcommand*\listfigurename{List of Figures}
\else
  \newcommand\listfigurename{List of Figures}
\fi
\ifdefined\listtablename
  \renewcommand*\listtablename{List of Tables}
\else
  \newcommand\listtablename{List of Tables}
\fi
\ifdefined\figurename
  \renewcommand*\figurename{Figure}
\else
  \newcommand\figurename{Figure}
\fi
\ifdefined\tablename
  \renewcommand*\tablename{Table}
\else
  \newcommand\tablename{Table}
\fi
}
\@ifpackageloaded{float}{}{\usepackage{float}}
\floatstyle{ruled}
\@ifundefined{c@chapter}{\newfloat{codelisting}{h}{lop}}{\newfloat{codelisting}{h}{lop}[chapter]}
\floatname{codelisting}{Listing}
\newcommand*\listoflistings{\listof{codelisting}{List of Listings}}
\makeatother
\makeatletter
\makeatother
\makeatletter
\@ifpackageloaded{caption}{}{\usepackage{caption}}
\@ifpackageloaded{subcaption}{}{\usepackage{subcaption}}
\makeatother
\usepackage{bookmark}
\IfFileExists{xurl.sty}{\usepackage{xurl}}{} % add URL line breaks if available
\urlstyle{same}
\hypersetup{
  pdftitle={A Mostly Harmless Introduction to Claude Code{,} Codex{,} and Friends},
  pdfauthor={Daniel Sánchez Pazmiño},
  pdflang={en},
  colorlinks=true,
  linkcolor={blue},
  filecolor={Maroon},
  citecolor={Blue},
  urlcolor={blue},
  pdfcreator={LaTeX via pandoc}}


\title{A Mostly Harmless Introduction to Claude Code, Codex, and
Friends}
\usepackage{etoolbox}
\makeatletter
\providecommand{\subtitle}[1]{% add subtitle to \maketitle
  \apptocmd{\@title}{\par {\large #1 \par}}{}{}
}
\makeatother
\subtitle{Study guide for empirical researchers}
\author{Daniel Sánchez Pazmiño}
\date{}
\begin{document}
\maketitle


\subsection{Overview}\label{overview}

This workshop introduces agentic AI for empirical research. The central
distinction is between a static chat model, which answers questions
inside a chat window, and an agent that can inspect a local project,
write files, run commands, and examine the results. The examples focus
on data cleaning, statistical code, legacy replication repositories,
literature review, reusable instructions, connectors, and local models.

The recording is a practical orientation rather than a complete course
in any one software package. The most useful habit it recommends is to
keep the researcher in charge: give the agent a bounded task, inspect
what it changed, run the analysis yourself, and verify the result
against the research question.

\begin{tcolorbox}[enhanced jigsaw, arc=.35mm, bottomrule=.15mm, bottomtitle=1mm, breakable, colback=white, colbacktitle=quarto-callout-note-color!10!white, colframe=quarto-callout-note-color-frame, coltitle=black, left=2mm, leftrule=.75mm, opacityback=0, opacitybacktitle=0.6, rightrule=.15mm, title=\textcolor{quarto-callout-note-color}{\faInfo}\hspace{0.5em}{Note}, titlerule=0mm, toprule=.15mm, toptitle=1mm]

The source transcript was generated automatically and lightly cleaned.
Product names and code terms have been normalized from context where
speech recognition was uncertain. Examples below are adapted from the
demonstration and are not presented as verbatim screen captures.

\end{tcolorbox}

\subsection{Learning objectives}\label{learning-objectives}

After studying this guide, you should be able to:

\begin{itemize}
\tightlist
\item
  distinguish a chat model from an agent with access to a project and
  terminal;
\item
  choose a sensible first workflow for an R, Python, Stata, or mixed
  research repository;
\item
  write bounded prompts that reduce unnecessary decisions and token use;
\item
  protect confidential data and recognize prompt injection as a security
  risk;
\item
  use Git, plans, permissions, and project instructions to make agent
  work recoverable;
\item
  apply an agent to data cleaning, code review, legacy code,
  replication, and literature search;
\item
  understand when an MCP connector, a reusable skill, or a local model
  is appropriate.
\end{itemize}

\subsection{Prerequisites and scope}\label{prerequisites-and-scope}

The talk assumes basic familiarity with files, a terminal, and at least
one research tool such as R, RStudio, Python, Stata, or MATLAB. A local
agent is most useful when the corresponding software is installed,
because the agent uses the computer's own programs to execute code. Some
hosted agent features also require a paid subscription.

The examples use a local folder containing compressed Ecuadorian
employment-survey files, R scripts, a legacy replication project, a
paper in PDF form, and supporting research material. The exact project
is less important than the pattern: an agent works best when the source
files, instructions, scripts, and outputs are visible in one bounded
project.

\subsection{1. From static chat to an
agent}\label{from-static-chat-to-an-agent}

A static model can explain a function, summarize text, or suggest code.
It is constrained by what is pasted or uploaded to the chat and usually
cannot inspect the actual state of a local research project. An agent
adds a loop around the model:

\begin{enumerate}
\def\labelenumi{\arabic{enumi}.}
\tightlist
\item
  inspect files and available tools;
\item
  propose or write a change;
\item
  run a command or statistical program;
\item
  inspect the output;
\item
  revise the work when the result does not match the request.
\end{enumerate}

That loop is the important change. An agent can read raw data, create a
cleaning script, call R from a terminal, and inspect the resulting CSV.
It can also organize folders or explain an old repository. A chat model
may describe those steps, but it does not automatically perform them on
the researcher's computer.

The talk uses a cake analogy for token limits. A subscription provides a
limited allocation of tokens, and more complex prompts and longer files
consume more of that allocation. Stronger models can be useful for
planning or difficult review, while a faster model is often enough for
routine implementation. The practical lesson is to spend context
carefully: give the agent the relevant files and a precise task instead
of asking it to make every decision from scratch.

An adapted version of the workshop's first project prompt might look
like this:

\begin{Shaded}
\begin{Highlighting}[]
\NormalTok{Inspect the files in case\_one. They are monthly compressed employment{-}survey files.}
\NormalTok{Create a clean dataset with one row per Ecuadorian province and one column per month.}
\NormalTok{Calculate median labour income using the survey weights. Preserve the raw files,}
\NormalTok{explain every transformation, and tell me which commands you ran.}
\end{Highlighting}
\end{Shaded}

The prompt defines the input, the unit of observation, the outcome, and
a safety constraint. It still leaves room for the agent to inspect the
data dictionary and ask about an ambiguity, but it does not ask for an
unbounded research project.

\subsection{2. Security, confidentiality, and
recoverability}\label{security-confidentiality-and-recoverability}

An online agent may work with local files, but information it reads can
still be sent to the provider. The talk therefore advises against using
online agents with social-security numbers, health identifiers,
personally identifiable information, or data protected by a data-sharing
agreement. A local open-source model can keep data on the researcher's
computer, but it requires suitable hardware and a more involved setup.

Prompt injection is a second risk. Untrusted web pages, documents, or
copied text can contain instructions aimed at the agent rather than
information for the researcher. A request to summarize a page can
therefore carry hidden instructions to open files, disclose information,
or run a command. Treat external content as data, not as an instruction
source, and use only sources you trust.

Agents can also make ordinary mistakes with serious consequences. A
mistaken command can overwrite files, delete a folder, or apply a
transformation to the wrong data. The talk recommends several layers of
protection:

\begin{itemize}
\tightlist
\item
  use Git or another recoverable backup for code and important project
  files;
\item
  work in a bounded project folder and use a sandbox when available;
\item
  select a plan or manual permission mode for unfamiliar tasks;
\item
  ask the agent to explain what it did and preserve raw inputs;
\item
  run tests and inspect the code before trusting a result.
\end{itemize}

\begin{tcolorbox}[enhanced jigsaw, arc=.35mm, bottomrule=.15mm, bottomtitle=1mm, breakable, colback=white, colbacktitle=quarto-callout-warning-color!10!white, colframe=quarto-callout-warning-color-frame, coltitle=black, left=2mm, leftrule=.75mm, opacityback=0, opacitybacktitle=0.6, rightrule=.15mm, title=\textcolor{quarto-callout-warning-color}{\faExclamationTriangle}\hspace{0.5em}{Warning}, titlerule=0mm, toprule=.15mm, toptitle=1mm]

An agent's confidence is not evidence that an analysis is correct. The
researcher remains responsible for the specification, identification
strategy, variable definitions, statistical assumptions, and final
interpretation.

\end{tcolorbox}

\subsection{3. A safe first workflow in a project
folder}\label{a-safe-first-workflow-in-a-project-folder}

The demonstration opens a local folder in Claude Code, points the agent
to the workshop repository, and asks it to inspect the first case. The
same pattern applies to Codex, Cursor, and GitHub Copilot: open the
project, state the task, let the agent inspect before editing, and keep
the raw data separate from derived outputs.

The agent can use a graphical application or a terminal. Working from a
terminal inside RStudio is useful when the researcher wants to keep the
familiar IDE in view. The command shown in the talk is conceptually
equivalent to starting the agent from the current project directory:

\begin{Shaded}
\begin{Highlighting}[]
\FunctionTok{cd}\NormalTok{ path\textbackslash{}to\textbackslash{}research{-}project}
\NormalTok{claude}
\end{Highlighting}
\end{Shaded}

The exact command depends on the installed CLI. The important property
is the working directory: the agent should start inside the repository
it is meant to inspect.

Three operating modes are especially useful:

\begin{itemize}
\tightlist
\item
  \textbf{Automatic mode:} convenient for a routine, well-bounded task
  when the project is backed up.
\item
  \textbf{Manual mode:} the agent asks before running commands or making
  changes.
\item
  \textbf{Plan mode:} the agent writes a proposed sequence without
  implementing it, which is useful for a large or uncertain change.
\end{itemize}

The recording also discusses memory. A provider may retain preferences
from previous conversations if memory is enabled. This can make later
work more consistent, but it can also preserve a mistaken assumption.
Inspect, correct, or disable memory when the project requires a clean
context.

Project instructions are a more explicit and portable control. Claude
Code reads \texttt{CLAUDE.md}; other agents commonly read
\texttt{AGENTS.md}. A small file at the repository root can say, for
example:

\begin{Shaded}
\begin{Highlighting}[]
\FunctionTok{\# Project instructions}

\SpecialStringTok{{-} }\NormalTok{Write readable, script{-}first statistical code.}
\SpecialStringTok{{-} }\NormalTok{Preserve raw data and explain derived files.}
\SpecialStringTok{{-} }\NormalTok{Run the relevant checks before reporting completion.}
\SpecialStringTok{{-} }\NormalTok{Ask before changing the research specification.}
\SpecialStringTok{{-} }\NormalTok{Use the project\textquotesingle{}s preferred language and naming conventions.}
\end{Highlighting}
\end{Shaded}

The talk recommends keeping these instructions close to the project and
updating them as new model behaviour creates new problems. They are not
a substitute for review, but they reduce repeated corrections.

For code, a minimal recoverability cycle is:

\begin{Shaded}
\begin{Highlighting}[]
\FunctionTok{git}\NormalTok{ status}
\FunctionTok{git}\NormalTok{ add path/to/intended/files}
\FunctionTok{git}\NormalTok{ commit }\AttributeTok{{-}m} \StringTok{"Save validated analysis change"}
\end{Highlighting}
\end{Shaded}

These commands are a compact illustration of the principle discussed in
the talk: create a checkpoint before allowing an agent to make a
substantial change, and keep the remote repository available as a
recovery path.

\subsection{4. Data cleaning and statistical
code}\label{data-cleaning-and-statistical-code}

The first live example begins with monthly compressed survey files from
Ecuador. The agent unzips and organizes them, derives a province
variable, identifies the income field, notices special or top-coded
values, and writes R code to produce a province-by-month dataset. It
also computes a survey-weighted median of labour income.

This is valuable because the mechanical work is time-consuming, but the
example also shows the limits of automation. The agent made decisions
the speaker had not explicitly requested, including details of the
weighted calculation. That can be helpful when the choice is correct and
dangerous when it is not. The researcher should inspect the data
dictionary, the handling of missing and special values, the survey
design, and the final sample before using the result.

The recommended division of labour is:

\begin{enumerate}
\def\labelenumi{\arabic{enumi}.}
\tightlist
\item
  ask the agent to inspect and propose a transformation;
\item
  allow it to write a readable script and a short explanation;
\item
  open the script in RStudio or the relevant IDE;
\item
  run it in the actual research environment;
\item
  compare diagnostics and outputs with the intended design.
\end{enumerate}

The agent can run code through the terminal when R or another program is
installed. It is not executing the analysis inside a magical cloud
version of R. It is calling the researcher's local installation and then
reading the result. This is why a missing R, Stata, or Python
installation prevents the workflow from completing.

The talk notes that language models are trained heavily on
software-engineering code. They may therefore produce abstractions,
loops, or modules that are technically valid but unnecessarily hard for
a researcher to read. Readability is not cosmetic in empirical work: it
supports review, replication, and later correction. State the preferred
style in the prompt or in project instructions.

\subsection{5. Legacy repositories, replication, and
translation}\label{legacy-repositories-replication-and-translation}

An agent is often most useful when the code already exists. The speaker
points it to an old replication folder and asks for an explanation
without running or modifying anything. This is a good first step for a
repository that has not been opened for several years, or for a
replication package downloaded from a journal.

The agent can summarize the scripts, identify the data flow, locate
graphs and side scripts, and explain how the final results are produced.
It can also propose a translation from R or Python to Stata. Running
Stata requires a valid local installation and, in the talk's workflow,
an additional Stata MCP connector. The connection is described as less
seamless than the one for open-source languages, partly because there is
less public training material for some closed-source software.

An adapted review prompt is:

\begin{Shaded}
\begin{Highlighting}[]
\NormalTok{Review this legacy replication repository. Do not modify files and do not run code.}
\NormalTok{Explain the data flow, the main scripts, the outputs, and any assumptions that}
\NormalTok{would need to be checked before reproducing the reported results.}
\end{Highlighting}
\end{Shaded}

For a real replication audit, the agent can go further: run the code,
compare the reported tables with the generated output, locate
inconsistencies, and propose a correction. The talk presents this as a
way to make code review and replication checks less expensive. It still
requires an expert to verify whether the code implements the paper's
intended specification and whether the data and environment are
comparable.

\subsection{6. Literature review and document
formats}\label{literature-review-and-document-formats}

The workshop compares three common inputs for a literature task. A
Markdown file is the easiest for an agent to read and costs less context
than a PDF. A Word document is usually easier than a PDF. PDFs can be
read, but layout, columns, tables, and scanned text make extraction more
expensive and less reliable.

For a paper already stored in a project, a sensible sequence is:

\begin{enumerate}
\def\labelenumi{\arabic{enumi}.}
\tightlist
\item
  convert the document to a reviewable text format when confidentiality
  permits;
\item
  ask the agent to summarize the argument and identify the literature
  gap;
\item
  ask it to search for additional papers or use a paper-search
  connector;
\item
  inspect the papers and the proposed citations yourself;
\item
  write only the claims that the sources actually support.
\end{enumerate}

The talk mentions connectors for PubMed, Google Scholar, and other
paper-search services. Some require an account or an institutional
subscription. Install a connector only from a trusted source and limit
it to the current project when the agent offers that choice.

An adapted prompt for this workflow is:

\begin{Shaded}
\begin{Highlighting}[]
\NormalTok{Read the paper text in case\_three. Summarize the literature review, identify}
\NormalTok{unsupported or missing claims, and propose three papers to investigate. Do not}
\NormalTok{rewrite the paper yet. Give the reason for each recommendation and distinguish}
\NormalTok{what you verified from what needs my review.}
\end{Highlighting}
\end{Shaded}

The goal is not to delegate the scholarly judgment. The agent can widen
the search and reduce mechanical effort, but it cannot decide whether a
contribution is genuinely new or whether a research design is credible.

\subsection{7. Reusable instructions, skills, and
MCPs}\label{reusable-instructions-skills-and-mcps}

The talk separates three kinds of customization:

\begin{itemize}
\tightlist
\item
  \textbf{Project instructions:} files such as \texttt{CLAUDE.md} or
  \texttt{AGENTS.md} that apply within a repository.
\item
  \textbf{Skills:} reusable Markdown instructions for a task such as
  proofreading, reviewing a paper, humanizing text, creating a lecture,
  checking a package, or reproducing an author's voice.
\item
  \textbf{MCPs:} connectors that give an agent an additional capability
  or a structured way to access a source, such as a paper database, a
  government data service, or Stata.
\end{itemize}

A skill is best understood as a reusable process description rather than
a model or a program. It can specify the inputs, the sequence of checks,
the output format, and the quality standard. The speaker describes
building a writing-voice skill from several papers and then applying it
to later drafts. The same pattern can be used for code conventions or a
replication workflow.

The maintenance cost matters. New model releases can introduce new
habits, such as overly abstract code or an unnatural writing style.
Treat instructions as living project assets: test them on a small task,
review the output, and update them when a recurring failure appears.
Keep project-specific rules in the project and general rules in a
reusable skill.

MCPs can extend an agent beyond its default tools, but they also expand
the trust boundary. A connector that reads email, queries a data
service, or runs Stata should be enabled only when needed, scoped to the
smallest useful project, and reviewed like any other software
dependency.

\subsection{8. Local models and sensitive
research}\label{local-models-and-sensitive-research}

For restricted data, the talk recommends considering a local open-source
model rather than sending information to an online provider. A
university or employer may host a model for researchers so that data
remains within its infrastructure. Individuals can also run smaller
models, but the tradeoff is hardware: ordinary computers cannot run the
most capable models at comfortable speed.

The key distinction is not that a local model is automatically correct
or safe. It changes where the data is processed. You still need access
controls, backups, review, and a clear understanding of what software is
running. A local model may also be weaker at niche statistical code, so
the researcher should test it on representative tasks before relying on
it.

\subsection{Cheat sheet}\label{cheat-sheet}

{\def\LTcaptype{none} % do not increment counter
\begin{longtable}[]{@{}
  >{\raggedright\arraybackslash}p{(\linewidth - 4\tabcolsep) * \real{0.3333}}
  >{\raggedright\arraybackslash}p{(\linewidth - 4\tabcolsep) * \real{0.3333}}
  >{\raggedright\arraybackslash}p{(\linewidth - 4\tabcolsep) * \real{0.3333}}@{}}
\toprule\noalign{}
\begin{minipage}[b]{\linewidth}\raggedright
Need
\end{minipage} & \begin{minipage}[b]{\linewidth}\raggedright
Agent pattern
\end{minipage} & \begin{minipage}[b]{\linewidth}\raggedright
Minimal example
\end{minipage} \\
\midrule\noalign{}
\endhead
\bottomrule\noalign{}
\endlastfoot
Inspect a project & Start in the project root and ask for an inventory &
\texttt{Inspect\ the\ repository.\ Do\ not\ modify\ files.} \\
Bound a task & State inputs, outputs, constraints, and checks &
\texttt{Preserve\ raw\ files\ and\ report\ every\ transformation.} \\
Protect changes & Use a checkpoint before a substantial edit &
\texttt{git\ status} then a focused commit \\
Control execution & Choose automatic, manual, or plan mode &
\texttt{Plan\ the\ changes;\ do\ not\ implement\ yet.} \\
Set preferences & Add repository instructions & \texttt{CLAUDE.md} or
\texttt{AGENTS.md} \\
Clean data & Ask for a script and diagnostics, then run it yourself &
\texttt{Write\ the\ cleaning\ script\ and\ explain\ special\ values.} \\
Understand legacy code & Request an explanation without edits first &
\texttt{Explain\ the\ data\ flow;\ do\ not\ run\ or\ modify\ code.} \\
Review replication & Compare paper claims, code, data, and outputs &
\texttt{Identify\ mismatches\ and\ show\ how\ to\ verify\ them.} \\
Search literature & Use text input and a scoped paper connector &
\texttt{Propose\ papers\ and\ separate\ verified\ claims\ from\ leads.} \\
Access structured data & Use a trusted data MCP &
\texttt{Query\ the\ official\ source\ and\ save\ the\ raw\ response.} \\
Protect restricted data & Use an approved local or institutional model &
\texttt{Do\ not\ send\ this\ dataset\ to\ an\ external\ service.} \\
\end{longtable}
}

\subsection{Final takeaways}\label{final-takeaways}

Agentic AI is most useful in research when it removes mechanical
friction without removing human control. Give it a bounded project,
explicit instructions, and a recoverable working state. Let it inspect
files, write scripts, run routine commands, explain legacy code, search
sources, and produce drafts. Then verify the data, code, specification,
citations, and interpretation yourself.

The workshop's practical sequence is therefore: protect the project,
state the task precisely, inspect the agent's plan, let it work within a
defined boundary, run and review the result, and turn repeated
preferences into project instructions or reusable skills. For
confidential data, move the workflow to an approved local or
institutional model rather than assuming that a local folder alone
prevents external processing.

\subsection{Resources named in the
recording}\label{resources-named-in-the-recording}

The recording discusses Claude Code, Codex, Cursor, GitHub Copilot,
RStudio, Git, Stata MCP, paper-search MCPs, PubMed, Google Scholar,
government-data MCPs, local open-source models, \texttt{CLAUDE.md}, and
\texttt{AGENTS.md}. These names are retained as the workshop's resource
map; no additional links have been added beyond what appears in the
recording.




\end{document}
